\setcounter{section}{0}
\section{Elektrische Schwingungen und Funkwellen}

\subsection{Zehnerpotenzen}

\begin{sol}{EA108}
 Rechnung: $$ 420 \cdot 10^{-6}\si{\ampere}  = \SI{0,00042}{\ampere}$$
\end{sol}

\begin{sol}{EA109}
 Rechnung: $$ 42 \cdot 10^{-3}\si{\ampere}  = \SI{0,042}{\ampere}$$
\end{sol}

\begin{sol}{EA110}
 Rechnung: $$ 4,2 \cdot 10^{6}\si{\hertz}  = \SI{4200000}{\hertz}$$
\end{sol}

\begin{sol}{EA111}
 Rechnung: $$ 10 \cdot 10^{-6}\si{\volt}  = \SI{0,00001}{\volt} =\SI{0,01}{\milli\volt}$$
\end{sol}

\begin{sol}{EA112}
 Rechnung: $$ 2 \cdot 10^{3}\si{\ohm}  = \SI{2000}{\ohm} =\SI{0,002}{\mega\ohm}$$
\end{sol}

\begin{sol}{EA113}
 Rechnung: $$ 2 \cdot 10^{-7}\si{\watt}  = \SI{0.0000002}{\watt} =\SI{0,2}{\micro\watt}$$
\end{sol}

\begin{sol}{EA114}
 Rechnung: $$ 5 \cdot 10^{-1}\si{\watt}  = \SI{0.5}{\watt} =\SI{500}{\milli\watt}$$
\end{sol}

\begin{sol}{EA115}
  $$ \SI{0,22}{\micro\farad}  = \SI{220}{\nano\farad} $$
\end{sol}

\begin{sol}{EA116}
  $$ \SI{3750}{\kilo\hertz}  = \SI{3,750}{\mega\hertz} $$
\end{sol}

\begin{ohmchapter}
Am Anfang der Formelsammlung findest Du eine Tabelle mit Zehnerpotenzen. Es lohnt sich aber dennoch zu üben, die Umrechnung von Zehnerpotenz brauchst Du später an vielen Stellen des Stoffs der Klasse E.
\end{ohmchapter} 

\subsection{Wellenlänge II}

\begin{sol}{EB311}
  Der erfahrene Funkamateure weiss, dass \SI{1.84}{\mega\hertz} im \SI{160}{\meter} Band liegt.
  Wir können aber auch rechnen:
  $$\lambda[m] = \frac{300}{f[MHz]} = \frac{300}{\SI{1.84}{\mega\hertz}}\approx \SI{163}{\meter} $$
\end{sol}

\begin{sol}{EB312}
  Rechnung:
  $$\lambda[m] = \frac{300}{f[MHz]} = \frac{300}{\SI{21}{\mega\hertz}}\approx \SI{14,29}{\meter} $$
\end{sol}

\begin{sol}{EB313}
  Rechnung:
  $$\lambda[m] = \frac{300}{f[MHz]} = \frac{300}{\SI{28.5}{\mega\hertz}}\approx \SI{10,5}{\meter} $$
\end{sol}

\begin{sol}{EB314}
  Rechnung:
  $$ f[MHz] = \frac{300}{\lambda[m]} = \frac{300}{\SI{80}{\meter}}\approx \SI{3.75}{\mega\hertz} $$
\end{sol}

\begin{sol}{EB315}
  Rechnung:
  $$ f[MHz] = \frac{300}{\lambda[m]} = \frac{300}{\SI{0,03}{\meter}}= \SI{10}{\giga\hertz} $$
\end{sol}

\begin{sol}{EB316}
  Rechnung:
  $$ f[MHz] = \frac{300}{\lambda[m]} = \frac{300}{\SI{0,1}{\meter}}= \SI{3}{\giga\hertz} $$
\end{sol}



\begin{ohmchapter}
    Es gelten die Formeln:
     $$ f[MHz] = \frac{300}{\lambda[m]} $$
     und $$\lambda[m] = \frac{300}{f[MHz]} $$

     Vereinfacht kannst Du Dir merken 
     $$ \text{gesuchter Wert} = \frac{300}{\text{gegebener Wert}}.$$
\end{ohmchapter} 
